% Sources: PAPER_MASTER_PLAN sec 9.2 (audit 2026-06-19) + sec 10.1 re-sweep
% (2026-06-26). Verify TODO-VERIFY bib entries before submission.

\paragraph{Security of collaborative perception.}
Attacks on, and defenses of, shared perception have been studied intensively
for connected autonomous vehicles. Zhang et al.~\cite{zhang2024cad}
demonstrate practical data-fabrication attacks on vehicular collaborative
perception and propose CAD, a detector in which benign vehicles jointly
reveal malicious fabrication through cross-vehicle occupancy-consistency
checks; the
approach therefore relies on benign vehicles positioned to observe the
attacked region. Two aspects of our setting go beyond this. First, we do not
assume such a benign observer of the attacked region: our test is pairwise
and remains effective with up to three traitors among ten robots. Second, we
specifically target \emph{camouflaged} fabrications --- phantoms blended into
real geometry --- which are the hardest to separate from honest observation
error and which a single-frame cross-agent check cannot contradict under
sensor noise. In feature-fusion collaborative perception,
GCP~\cite{gcp2025} defends against malicious agents by augmenting
single-shot spatial-consistency checks with temporal motion-flow
reconstruction to expose blind-area attacks; MADE~\cite{made2024} detects
and removes malicious agents through hypothesis tests on the consistency
between each inspected agent and the ego agent, with statistical
false-positive control; and CoDynTrust~\cite{codyntrust2025} addresses the
benign counterpart by weighting intermediate feature contributions
according to their estimated uncertainty, thereby mitigating temporal
asynchrony rather than adversarial behaviour; and ROBOSAC~\cite{robosac2023}
rejects adversarial feature-map perturbations by sampling subsets of teammates
until the collaborative output reaches consensus with the ego's own perception,
falling back to ego-only perception when no consensus is found. All four are
evaluated using Average Precision (AP) on vehicular benchmarks.

A second, closer line verifies \emph{object-level claims against the
verifier's own observations}. TruPercept~\cite{trupercept2020} evaluates reported
object detections as verified locally, weighting each evaluation by the
visibility of the claimed object from the verifier's viewpoint before
computing peer trust; MATE~\cite{mate2025} formalizes this as Bayesian trust
estimation over geometric observation consistency with explicit
field-of-view reasoning, against false-positive, false-negative, and
translation attacks on multi-target tracking; and the same authors' aerial
framework~\cite{aerialtrust2025} carries that hidden-Markov trust estimation
to ad hoc UAV networks on intelligence, surveillance, and reconnaissance
(ISR) missions --- the setting nearest to ours in the literature. Our single-frame filters (Section~\ref{sec:methods-defense})
belong to precisely this family, and inherit its central assumption: that a
verifier positioned to observe a claimed object either corroborates or
contradicts it.

That assumption is what a camouflaged fabrication attacks. It is instructive
that TruPercept additionally rejects \emph{open-space} phantoms with a
plausibility check that culls the sensor frustum toward a claimed object: an
object with no returns in front of it, but returns behind it, is deemed
absent with high confidence~\cite{trupercept2020}. Operationally, such a
test derives its evidence from rays passing through empty space at the
claimed location, which is exactly what a phantom placed in free space provides ---
and, in our own experiments, one reason the conspicuous wall attack proves
comparatively easy to detect (Section~\ref{sec:res-temporal}). A camouflaged
phantom is designed to avoid supplying that evidence: placed flush against a
real obstacle, its claimed volume coincides with returns the verifier
genuinely receives from nearby structure, so a free-space plausibility test
has little to discriminate on. We note that TruPercept does not evaluate
camouflaged fabrications, and we make no claim about the behaviour of its
checker in our setting; the point is that geometric plausibility tests
derive their power from rays traversing empty space, and camouflage attacks
intentionally deny them that. Under ranging noise the residual discrepancy
further shrinks below the tolerance a verifier must grant honest neighbours
(Section~\ref{sec:res-robust}). Our temporal test therefore does not seek a
stronger \emph{spatial} contradiction; it accumulates a \emph{statistical}
one over time.

This line accommodates natural sensor error through longitudinal filtering,
and observes that an attacker must remain dynamics-consistent to evade local
filtering; TruPercept further reports that unreliable agents which mostly
report truthfully are not separated from honest ones by its trust
model~\cite{trupercept2020} --- an early indication of the evasion regime we
formalize. Where our setting departs further:
TruPercept~\cite{trupercept2020} aggregates trust at a central server and
MATE~\cite{mate2025} at a central computing centre (the aerial
framework~\cite{aerialtrust2025} does distribute the estimation), whereas
our verdicts are pairwise and local, as a swarm without infrastructure
requires. None of this line
characterizes the regime in which the consistency check \emph{itself} turns
harmful under ranging noise (our destructive-filter result), considers
fabrications geometrically camouflaged within the noise tolerance of real
objects, or measures impact on a closed-loop learned \emph{navigation} task
--- their trust protects a shared perception picture, whether for driving, a
smart city, or aerial surveillance, whereas ours protects the obstacle map a
policy steers by. Our contribution begins where this family breaks. Closest to our mechanism,
the recent PRBI defense~\cite{prbi2026} exploits frame-to-frame perceptual
consistency to identify lying vehicles in fully untrusted cooperative
detection. It differs from our setting in three ways that matter: it is
evaluated by detection accuracy on feature-level fusion rather than by a
closed-loop navigation objective; it does not model ranging noise, and so
does not confront the honest-disagreement regime in which we show
consistency checking becomes destructive; and while it evaluates
intermittent injection schedules, it does not consider a defense-aware
attacker that optimizes its fabrication against the deployed detector, as
our stealth/harm-bind analysis does. More fundamentally, PRBI's alarm is
raised when the Jaccard similarity between the current and preceding
detection sets falls below a threshold~\cite{prbi2026}: it screens for a
\emph{drop} in inter-frame agreement. Such a criterion responds to
perturbations of the scene relative to its recent past. A
persistent fabrication is temporally self-consistent once established;
PRBI does not evaluate this attack class, and we make no claim about how
its detector would behave against it. The
distinction we draw is one of reference signal: PRBI compares the scene to
its own past, whereas our test compares a neighbour's claim to the
verifier's independent sensing of the same object, accumulated over time.

Complementing the defense literature, the TrustFlip
attack~\cite{trustflip2026} shows that consistency-based trust can itself be
weaponized: physically induced disagreements cause defenses to expel
\emph{honest} vehicles. Our evaluation addresses the same failure class
by measuring \emph{no-harm} --- the success cost of running the defense in
attack-free swarms --- directly, and finding it statistically
indistinguishable from zero even under an adaptive attacker.

\paragraph{Temporal detection of sensor spoofing.}
A separate line detects LiDAR spoofing against a \emph{single} vehicle by
exploiting temporal structure in its own sensor stream, e.g.\ motion-induced
consistency in 3D-TC2~\cite{tc2_2021} and point-level temporal consistency
in ADoPT~\cite{adopt2023}. These methods test whether observations from a
single sensor remain temporally self-consistent. Ours instead aggregates a \emph{cross-agent
disagreement} --- the offset between a neighbour's claim and the verifier's
own view --- whose statistics under honest noise (zero-mean, variance
$2\sigma^2$) versus fabrication (persistent bias) we can characterize and
threshold explicitly. The statistic itself is elementary; the contribution
is its use as a communication-trust discriminator that survives sensor
noise and camouflage, conditions under which we show single-frame
cross-agent checks fail.

\paragraph{Byzantine resilience in multi-robot systems.}
Byzantine fault tolerance originates in distributed
computing~\cite{lamport1982byzantine} and, in robot swarms, is typically
addressed at the \emph{consensus} layer. Consensus-based coordination such
as SwarmRaft~\cite{swarmraft2025} fuses peer measurements to maintain
agreement on state (e.g.\ position and heading) under degraded conditions,
and evolutionary-game analyses~\cite{conformity2026} model how deceptive
strategies propagate through conformity dynamics under Byzantine influence.
These works protect state agreement or model attacks on collective
decision-making; our
threat lives one layer below, in the \emph{perception content} that robots
exchange, and our defense verifies that content against the verifier's own
sensing rather than against majority agreement. Because each verdict is
pairwise, the mechanism is structurally independent of the number of liars;
we confirm this empirically (Section~\ref{sec:res-temporal}), where its
recovery remains significant even when traitors form a local majority among
an honest robot's neighbours --- up to seven of ten robots. Sampling-based
consensus defenses such as ROBOSAC~\cite{robosac2023}, by contrast, must draw
an attacker-free subset of teammates and know or estimate the attacker ratio,
with a sampling budget that grows sharply as benign agents become scarce; our
pairwise test needs neither a clean subset, a known ratio, nor a consensus
vote.

\paragraph{Learned multi-robot navigation and robust communication.}
Our navigator is trained with centralized-training/decentralized-execution
multi-agent reinforcement learning, following the MAPPO
recipe~\cite{yu2022mappo} built on PPO~\cite{schulman2017ppo} and, in
implementation, on Stable-Baselines3~\cite{raffin2021sb3}; CTDE itself
traces to multi-agent actor-critic formulations such as
MADDPG~\cite{lowe2017maddpg}. Our attack surface is deliberately
interpretable: the messages are geometric obstacle claims rather than
learned latent vectors, the fusion rule is explicit, and the defense is a
hand-designed statistical test --- one of our findings is that no learned
trust component is required to counter this threat class.

The closed-loop learned cooperative-perception paradigm we build on --- a
policy that consumes shared perception and is scored by end-task success ---
was established for benign networked driving by
Coopernaut~\cite{coopernaut2022}; our contribution is not that paradigm but
the adversarial layer within it. Closest in driving-level outcome, the recent
PosePert attack~\cite{stealthyfab2026} propagates small, stealthy pose
perturbations of \emph{existing} objects through a modular tracking-and-
prediction stack to induce unsafe behaviour on a single vehicle, measured by
trajectory-prediction error and a scripted braking-risk rate. Our setting
differs in the attack (fabricated obstacles, not pose shifts of real objects),
the control layer (a learned navigation policy executing in closed loop, rather
than a tracker-plus-predictor pipeline), the end metric (task success, not
prediction error or a safety flag), and the multi-robot rather than
single-vehicle scope.

In combination, to our knowledge no prior work evaluates
fabricated-obstacle attacks on collaborative perception \emph{inside a
closed-loop learned navigation task} (success, not detection accuracy, as
the end metric), analyzes the interaction of consistency-based trust with
sensor noise to the point of demonstrating harmfulness, closes the
resulting camouflage blind spot with a cross-agent temporal bias test, and
stress-tests the closure against an adaptive attacker across noise levels
and traitor counts.

\paragraph{On baseline comparison.}
The defenses above cannot be transplanted onto our benchmark without
distortion: CAD, PRBI, GCP, and MADE defend the fused perception pipelines
of vehicular detection stacks and are evaluated by detection-level metrics
(anomaly-detection rates or average precision), while TruPercept, MATE, and
the same authors' aerial framework operate at the object and track level;
none is designed to operate directly on the geometric obstacle claims
exchanged by our agents, and none produces the navigation-success outcome
we measure --- a reimplementation would test our translation of each method
rather than the method itself. We
instead compare against the \emph{primitive that the object-level family
shares}: verifying a peer's claims against what the verifier itself was
positioned to observe, accumulating trust, and excluding on contradiction.
That primitive is implemented natively in our setting in two strengths ---
the fixed-tolerance (naive) variant and the noise-aware robust variant
(Section~\ref{sec:methods-defense}) --- and it is the same
visibility-and-consistency recipe underlying
TruPercept~\cite{trupercept2020} and the MATE
line~\cite{mate2025,aerialtrust2025}. Our evaluation therefore reports how
the family's core mechanism behaves in this regime --- including its
noise-induced failure --- before showing what the temporal test adds on top
of it.
